\section{Cards E: Power, Energy, Area, Reliability}\label{grp:power}
\begin{multicols}{2}
\card{E1}{DVFS (dynamic voltage-frequency scaling)}{Dynamic}
\fx{Energy/power too high for battery or thermal budget.}
\fh{Lower $V$ and $f$ when full speed is not needed. Dynamic power $=\tfrac12CV^2f$; energy per switch $=\tfrac12CV^2$.}
\fu{Battery, thermal cap, bursty load, cores throttling.}
\fb{Lowering clock alone: it cuts power but \emph{not} energy per task (slide); lowering $V$ cuts energy via $V^2$.}
\fa{Constant full load or hard real-time (use fixed V/f: static, predictable).}
\fw{Slower at low settings; transition latency.}
\fn{Per-core voltage/clock domains; OS/firmware governor.}
\fq{If $V$ scales with $f$, power $\propto f^3$: $0.8\times$ clock $\Rightarrow\approx0.51\times$ power\bs.}
\fj{Because [load is bursty / battery-limited], DVFS lowers $V$ in idle phases, cutting energy per task by the $V^2$ term, which clock gating alone cannot.}

\card{E2}{Power gating, turning off idle units/cores}{Dynamic}
\fx{Static (leakage) power $=I_{static}\times V$, which grows with transistor count.}
\fh{Cut supply to idle blocks/cores; low-power DRAM states.}
\fu{Many cores often idle, PMD, large caches mostly unused.}
\fb{DVFS for idle blocks: a block at low V still leaks; gating removes it.}
\fw{Wake-up latency; state must be saved.}
\fn{Wake-up policy; save/restore of state.}
\fj{Idle [cores/units] still leak; gating them removes static power, which DVFS does not.}

\card{E3}{Many small cores (or big + small) instead of one huge core}{Static}
\fx{Power wall, area cost, ILP saturation.}
\fh{Spend transistors on several simpler cores; optionally a few big cores for serial parts and many small ones for parallel parts\bs.}
\fu{Threaded/parallel software, area or power limit, throughput goal.}
\fb{One wide core: ILP saturates at 3--6 issue and ILP hardware is energy-inefficient; big dies cost disproportionately more (yield). Two slower cores can beat one fast core on throughput per watt.}
\fa{Software is serial (Amdahl): then keep one strong core.}
\fw{Serial phases run slower on small cores; coherence, sync and shared-bandwidth costs grow with core count.}
\fn{Threads/tasks to run; memory bandwidth and coherence for all cores.}
\fq{Yield $\propto(1+\text{defects}\times\text{area})^{-N}$, $N=11.5$--15.5: bigger die $\Rightarrow$ much costlier.}
\fj{Because the workload has [N independent tasks] and the budget is [power/area], I use [k] small cores rather than one wide core whose extra issue slots would stay idle.}

\card{E4}{Dependability: redundancy, ECC}{Static}
\fx{Downtime and memory errors (servers, 24/7).}
\fh{Error-correcting codes on memory/caches; redundant components; fast repair.}
\fu{Servers, 24/7 services, large DRAM capacity.}
\fw{Extra bits, area; ECC harder with wide memory (partial writes need read-modify-write).}
\fq{Availability $=\text{MTTF}/(\text{MTTF}+\text{MTTR})$.}
\fj{For 24/7 service, availability $=$ MTTF/(MTTF+MTTR); ECC and redundancy raise MTTF and cut MTTR.}

\columnbreak
\begin{key}[Technology-trend arguments (quote to justify direction)]
\begin{itemize}
\item \textbf{Memory wall:} processor bandwidth improved 10{,}000--25{,}000$\times$ but memory only 300--1200$\times$; latency 30--80$\times$ vs.\ 6--8$\times$. Bandwidth is easy to buy, latency is not $\Rightarrow$ caches, banks, latency hiding.
\item \textbf{Wires do not scale} like transistors $\Rightarrow$ huge centralised structures (wide issue, big CDB, many ports) hurt the clock $\Rightarrow$ smaller replicated cores.
\item \textbf{Single-thread growth ended} ($\sim$2003); ILP cannot be pushed further $\Rightarrow$ explicit DLP/TLP/RLP.
\item \textbf{Energy per task} is the right metric, not power.
\item \textbf{Focus on the common case} (Amdahl): speed up what dominates time.
\end{itemize}
\end{key}

\begin{key}[Machine classes: what each optimises]
\textbf{PMD}: energy, real-time. \textbf{Desktop}: price-performance, single thread. \textbf{Server}: availability, scalability, throughput. \textbf{Cluster/WSC}: request-level parallelism, cost and energy per request. \textbf{Supercomputer}: FP and fast network. \textbf{Embedded}: price. Full reference designs: \pg{sec:classes}.
\end{key}
\end{multicols}
